% TOP_PROBLEM: {"id": 30006980, "title": "UCT Problem for Separable Nuclear Complex C*-Algebras", "category_id": 8, "category": {"id": 8, "name": "analysis", "display_name": "Analysis", "description": "Limits, continuity, calculus, and function theory.", "slug": "analysis", "order_index": 8, "created_at": "2026-07-31T15:26:25.671Z"}, "status": "open"}
% FIRST_POSTED: 2026-09-27
% Standalone research entry 203; batch 203--207.
% Original section material: Pasted text(3).txt, source edition 22 September 2026.
% Research additions: 27 September 2026.
% Compile twice with: lualatex 30006980.tex
% !TeX program = lualatex
% Compile twice: lualatex open_problems_500.tex
% All references and prose are embedded; no BibTeX database or external inputs.
\documentclass[11pt,a4paper]{article}
\usepackage[margin=24mm,headheight=14pt]{geometry}
\usepackage{fontspec}
\setmainfont{Latin Modern Roman}
\setsansfont{Latin Modern Sans}
\setmonofont{Latin Modern Mono}
\usepackage{amsmath,amssymb,mathtools}
\usepackage{unicode-math}
\setmathfont{Latin Modern Math}
\usepackage{microtype}
\usepackage{enumitem,xurl,needspace}
\usepackage[dvipsnames]{xcolor}
\usepackage{hyperref}
\hypersetup{unicode=true,colorlinks=true,linkcolor=MidnightBlue,urlcolor=MidnightBlue,
 pdftitle={500 Mathematical Problem Statements},pdfauthor={Source-annotated compilation},bookmarksnumbered=true}
\usepackage{bookmark}
\usepackage{tocloft}
\setlength{\cftsecnumwidth}{3em}
\cftsetpnumwidth{2.5em}
\cftsetrmarg{4em}
\usepackage{titlesec}
\titleformat{\section}{\Large\bfseries\raggedright}{\thesection}{0.7em}{}
\usepackage{fancyhdr}
\pagestyle{fancy}\fancyhf{}
\fancyhead[L]{\small\sffamily Mathematical problem statements}
\fancyhead[R]{\small Source edition 22}
\fancyfoot[C]{\thepage}
\setlength{\parindent}{0pt}
\setlength{\parskip}{5pt plus 1pt}
\setlength{\emergencystretch}{3em}
\setcounter{tocdepth}{1}
\setcounter{secnumdepth}{1}
\setlist[itemize]{leftmargin=3em,itemsep=5pt,topsep=4pt,parsep=0pt}
\allowdisplaybreaks
\newcommand{\N}{\mathbb{N}}
\newcommand{\Z}{\mathbb{Z}}
\newcommand{\Q}{\mathbb{Q}}
\newcommand{\R}{\mathbb{R}}
\newcommand{\C}{\mathbb{C}}
\newcommand{\F}{\mathbb{F}}
\newcommand{\A}{\mathbb{A}}
\newcommand{\PP}{\mathbb P}
\newcommand{\E}{\mathbb{E}}
\newcommand{\eps}{\varepsilon}
\newcommand{\dd}{\,\mathrm d}
\newcommand{\sref}[2]{\hyperlink{source:#1:#2}{[S#2]}}
\newcommand{\eref}[2]{\hyperlink{extra:#1:#2}{[E#2]}}
% Turn the next switch off for a shorter reading copy. The quotations preserve
% every exactTarget field, including qualifications omitted from a short summary.
\newif\ifshowcatalogue
\showcataloguetrue
\newcommand{\cataloguescope}[1]{\ifshowcatalogue
 \paragraph{Exact scope in the supplied catalogue.}
 \begin{quote}\small #1\end{quote}\fi}

% Standalone wrapper; no external inputs or bibliography files are required.
\fancyhead[L]{\small\sffamily Research notebook: section 30006980}
\fancyhead[R]{\small 27 September 2026}
\hypersetup{pdftitle={Research attempt 203},pdfauthor={Research notebook}}
\setlength{\parskip}{4pt plus 1pt}
\setlist[itemize]{before=\small\interlinepenalty10000,itemsep=3pt}
\begin{document}
\setcounter{section}{0}
% ================================================================
% problemId: 30006980
% Canonical problem ID: 30006980
% exactTarget SHA-256: 2a71526b635df8557fe6fb1acc2462f14c7f5131a81a264ff1d8c70684b53a7c
\clearpage
\section*{\texorpdfstring{UCT Problem for Separable Nuclear Complex C*-Algebras}{UCT Problem for Separable Nuclear Complex C*-Algebras}}\label{30006980}
\subsection{Definitions and mathematical statement}
For separable complex $C^*$-algebras $A,B$, let $K_i$ be operator $K$-theory and $KK_j(A,B)$ Kasparov's bivariant groups, graded modulo two. Nuclearity of $A$ means its identity admits point-norm approximations by completely positive contractions factoring through finite-dimensional matrix algebras. Define
\[
 H_j=\bigoplus_{i=0,1}\operatorname{Hom}_{\Z}(K_i(A),K_{i+j}(B)),\quad
 E_j=\bigoplus_{i=0,1}\operatorname{Ext}^1_{\Z}(K_i(A),K_{i+j+1}(B)).
\]
The problem asks whether every separable nuclear $A$ has the canonical exact sequence
$0\to E_j\to KK_j(A,B)\to H_j\to0$ for every separable $B$ and each $j$. The maps are the Rosenberg--Schochet maps, including the canonical kernel identification; an abstract short exact sequence with arbitrarily chosen maps is insufficient. Neither unitality nor a natural splitting is required.
\par\smallskip\textit{Formulation and concept references:} \sref{30006980}{1}.
\cataloguescope{Decide whether EVERY separable nuclear complex C*-algebra A, unital or nonunital, satisfies the Rosenberg-Schochet UCT: for EVERY separable complex C*-algebra B, unital or nonunital, and each j in Z/2, the canonical UCT sequence 0 -> E\_j(A,B) -> KK\_j(A,B) -> H\_j(A,B) -> 0 is exact. Here H\_j(A,B)=direct\_sum\_(i=0,1) Hom\_Z(K\_i(A),K\_(i+j)(B)) and E\_j(A,B)=direct\_sum\_(i=0,1) Ext\_Z\textasciicircum{}1(K\_i(A),K\_(i+j+1)(B)), with indices modulo2. The right map is the natural map on K-theory induced by KK classes; the left is the canonical inverse extension-class identification of the kernel, not an arbitrary chosen injection. No canonical or natural splitting is required. Equivalently, every such A belongs to the Rosenberg-Schochet bootstrap category of separable nuclear C*-algebras, or is KK-equivalent to a separable commutative complex C*-algebra.}
\subsection{Short English statement}
Does operator K-theory, together with its extension term, compute KK-theory through the canonical universal coefficient sequence for every separable nuclear complex C*-algebra?
% BEGIN RESEARCH ADDITION 203
\subsection{Research attempt}

\paragraph{Current frontier.}
The unrestricted UCT is still listed as Problem II in the May 2026 revision
of Schafhauser--Tikuisis--White; neither nuclearity nor the classification
results quoted there remove the UCT hypothesis.\eref{30006980}{1}
Barlak--Li prove the UCT for separable nuclear algebras with a Cartan
subalgebra, and more generally for the nuclear twisted groupoid algebras
in their Theorem 1.1.\sref{30006980}{1}
Kirchberg's reduction permits concentrating on unital Kirchberg algebras
with vanishing $K$-theory. Here ``Kirchberg'' includes separability,
nuclearity, simplicity, and pure infiniteness.\eref{30006980}{2}
Controlled $KK$-theory supplies an inverse-limit description with a
potentially nonzero $\lim^1$ term; simply discarding that term replaces
$KK$ by a quotient and loses information.\eref{30006980}{3}

\paragraph{Attack route.}
Three routes suggest different missing inputs. A Cartan construction would
need a regular maximal abelian subalgebra with faithful conditional
expectation, not merely a large commutative subalgebra.\sref{30006980}{1}
A finite-complexity decomposition would need compatible pieces and overlap
control of the kind used by Willett--Yu, rather than arbitrary nuclear
factorizations.\eref{30006980}{4}
The route pursued below instead asks whether a nuclear approximation can
be promoted to a \emph{genuine} $KK$-factorization whose error is nilpotent.
The algebraic correction is exact; the substantial question is whether the
required classes and nilpotence can be constructed. Two explicit tests
then examine the two most tempting shortcuts: converting completely
positive maps into $KK$-classes, and identifying classes from all their
finite-stage restrictions.

\paragraph{Attempt.}
\textbf{1. A finite algebraic correction, with all coefficient algebras retained.}
Write $\mathcal B$ for the bootstrap category in $KK$. We use its standard
closure under $KK$-retracts: a localizing triangulated category with
countable coproducts is closed under retracts, by the mapping telescope of
an idempotent. For separable nuclear algebras, membership in $\mathcal B$
is equivalent to the canonical UCT in the question, for all separable
coefficient algebras.\sref{30006980}{1}\eref{30006980}{1}

\textbf{Lemma 203.1 (nilpotent-error domination).}
Suppose $A$ is separable nuclear, $D\in\mathcal B$, and there are even
classes
\[
 x\in KK_0(A,D),\qquad y\in KK_0(D,A),\qquad
 e=1_A-x\otimes_Dy\in KK_0(A,A)
\]
such that $e^N=0$ for some finite $N\ge1$. Then $A$ satisfies the exact
canonical UCT stated above.

\emph{Proof.}
All powers and products in $KK_0(A,A)$ are Kasparov products. Set
\[
 q=1_A+e+\cdots+e^{N-1},\qquad y'=y\otimes_Aq.
\]
Associativity and the finite geometric-series identity give
\[
 x\otimes_Dy'=(1_A-e)q=1_A-e^N=1_A.
\]
Thus $A$ is a $KK$-retract of $D$, hence belongs to $\mathcal B$. Applying
the bootstrap UCT gives precisely its natural map on $K$-theory and its
canonical kernel identification, simultaneously for $j=0,1$ and every
separable $B$. No splitting is chosen or claimed to be natural.
This argument uses neither unitality of $A,B$ nor nuclearity of $B$.
\hfill$\square$

The useful feature is that an exact factorization of $1_A$ is unnecessary
\emph{at the construction stage}; a nilpotent error would suffice. The
lemma is elementary categorical algebra, not a claim of a new permanence
theorem. In particular, ``small in norm'' is not a substitute for
nilpotence in this generally non-normed group.

\textbf{2. Nuclear approximations do not yet supply those classes.}
There is a particularly transparent exact completely positive retraction.
Define unital completely positive maps
\[
 \phi:\C^2\longrightarrow M_3(\C),\quad
 \phi(a,b)=\operatorname{diag}\bigl(a,b,(a+b)/2\bigr),
 \qquad
 \psi(T)=(T_{11},T_{22}).
\]
The diagonal coordinates of $\phi$ are positive functionals, and the
coordinates of $\psi$ are vector states, so complete positivity follows
directly. Moreover $\psi\phi=\operatorname{id}_{\C^2}$. Nevertheless,
for the projection $p=(1,0)$,
\[
 \phi(p)-\phi(p)^2=\operatorname{diag}(0,0,1/4),\qquad
 \|\phi(p)-\phi(p)^2\|=1/4.
\]
Even an \emph{exact} completely positive retraction therefore does not force
its outward leg to be multiplicative. Repeating this pair at every stage
also refutes the assertion that improving the composite alone forces
asymptotic multiplicativity of each leg.

A second test keeps the topology visible. For increasingly fine finite
sets $\{z_{N,j}\}_{j=1}^{r_N}\subset S^1$, choose subordinate nonnegative
partitions of unity $\{h_{N,j}\}$. Put
\[
 \phi_N(f)=(f(z_{N,j}))_j,\qquad
 \psi_N((a_j))=\sum_j a_jh_{N,j}.
\]
These maps factor through $\C^{r_N}$, and uniform continuity gives
$\psi_N\phi_N(f)\to f$ in norm for every $f\in C(S^1)$.
Suppose, incorrectly, that one could regard both legs as even
$KK$-classes and obtain a product inducing the identity on $K$-theory.
The resulting map on $K_1(C(S^1))=\Z$ would factor through
$K_1(\C^{r_N})=0$, and hence be zero. This is a contradiction.
The example is not a counterexample to the UCT: $C(S^1)$ is already
commutative. It identifies exactly why completely positive approximation
alone cannot prove Lemma 203.1's hypothesis.

\textbf{3. A coherent-limit obstruction that survives every finite test.}
Let
\[
 A_n=C(S^1),\qquad \alpha_n(f)(z)=f(z^2),\qquad
 A=\varinjlim(A_n,\alpha_n).
\]
All connecting maps are injective. The limit is commutative and is in the
bootstrap category. The standard Milnor sequence, in the form recalled
in Willett--Yu, applies here; it is not being assumed for arbitrary
completely positive approximations.\eref{30006980}{3}
For $B=\C$ it gives
\[
 0\longrightarrow\varprojlim{}^{1} KK_1(A_n,\C)
 \longrightarrow KK_0(A,\C)
 \longrightarrow\varprojlim KK_0(A_n,\C)\longrightarrow0.
 \tag{203.1}
\]
The finite-stage groups are $KK_0(A_n,\C)=\Z$ and
$KK_1(A_n,\C)=\Z$. Pullback by $\alpha_n$ is the identity in degree zero
and multiplication by two in degree one. These assertions follow from
the degree of $z\mapsto z^2$ and the UCT for the circle, not from the UCT
being investigated for a general $A$.

Here the obstruction can be computed without an abstract vanishing
claim. For the inverse system $\Z\xleftarrow{2}\Z\xleftarrow{2}\cdots$,
index stages by $n\ge0$ and use the usual concrete definition
\[
 \varprojlim{}^{1}=\operatorname{coker}\Delta,\qquad
 \Delta:\prod_{n\ge0}\Z\longrightarrow\prod_{n\ge0}\Z,
 \quad\Delta(b)_n=b_n-2b_{n+1}.
\]
Let $\Z_2$ denote the $2$-adic integers. The convergent $2$-adic series
\[
 \Phi(a)=\sum_{n\ge0}2^n a_n\pmod{\Z}
\]
induces an isomorphism
\[
 \operatorname{coker}\Delta\ \cong\ \Z_2/\Z. \tag{203.2}
\]
Indeed, telescoping gives $\Phi(\Delta b)=b_0\pmod\Z=0$; binary digits
show surjectivity. If the series equals an ordinary integer $b_0$, define
\[
 b_n=2^{-n}\left(b_0-\sum_{i<n}2^ia_i\right).
\]
Its numerator is divisible by $2^n$ in $\Z$, so each $b_n$ is an integer,
and $b_n-2b_{n+1}=a_n$. This proves the kernel assertion.

The sequence $a=(1,0,1,0,\ldots)$ maps to
$1+4+16+\cdots=-1/3\in\Z_2$, which is not an ordinary integer. It
therefore represents a nonzero class in the left term of (203.1).
Its image in $KK_0(A,\C)$ restricts to zero on every $A_n$.
Consequently agreement on every finite stage does not imply equality in
$KK$. This example demonstrates a genuine nonzero \emph{phantom} class,
not a numerical failure to detect a class.

For comparison, if an inverse sequence has surjective bonding maps,
$\Delta$ is surjective: choose $b_0$, then solve recursively
$f_n(b_{n+1})=b_n-a_n$. This gives a valid mechanism for eliminating a
$\lim^1$ obstruction. In the displayed example multiplication by two is
not surjective and its images do not stabilize. A proposed construction
must supply such compatibility when it uses it; surjectivity cannot be
inferred from nuclearity. Nor should all extension information be
eliminated: here the nonzero term is precisely compatible with the UCT
extension contribution from $K_1(A)=\Z[1/2]$.

\textbf{4. Testing the nilpotence step against the minimal reduction.}
Consider the remaining Kirchberg class from the established reduction:
$A$ unital with $K_0(A)=K_1(A)=0$.\eref{30006980}{2}
The identity $1_A\in KK_0(A,A)$ is then invisible to $K$-theory. Suppose
one proposed the general rule that every such invisible endomorphism is
nilpotent. Applied to $1_A$, this says
\[
 0=(1_A)^N=1_A.
\]
This would imply $KK_j(A,B)=0$ for every $B,j$, since any class is its
product with $1_A$. Both $E_j$ and $H_j$ vanish, so the full canonical
UCT follows for this class, and the reduction then settles the original
problem. Thus this apparently technical ``nilpotence of the invisible
error'' is already as deep as the unrestricted UCT in the critical case.
It is not a consequence of $K_*(A)=0$ that may be used for free.

\paragraph{Stress test.}
The finite geometric-series argument has no convergence issue and retains
both parities and all separable coefficient algebras. The first failure
occurs earlier: no $KK$-factorization of the stipulated kind has been
obtained from a nuclear approximation. The exact $\C^2$ example rules out
one possible justification, while the circle example rules out replacing
its intermediate algebra by a finite-dimensional $KK$-model.

The solenoid calculation rules out a separate inference from finite-stage
agreement to global equality. It does not contradict controlled
$KK$-theory; it illustrates why its derived-limit term is needed.
Finally, invoking Kirchberg--Phillips classification \emph{with} the UCT
hypothesis to identify an arbitrary $K$-trivial Kirchberg algebra with
$\mathcal O_2$ would be circular. No such identification was used.

\Needspace{10\baselineskip}
\paragraph{Outcome.}
\textbf{Outcome: Conditional reduction.}

A genuine bootstrap factorization with nilpotent error is sufficient for
the exact UCT, and the correction is explicitly given. The notebook also
provides two completely positive countertests and computes a nonzero
finite-stage-invisible $KK$-class. These are diagnostic deductions, not a
claim that a previously untreated nuclear algebra has entered the
bootstrap category.

What remains is a construction of $D,x,y$ with the required nilpotence,
or another argument carrying the missing coherent extension information.
The $K$-trivial Kirchberg test shows that a blanket nilpotence assertion
would already resolve the full problem. The conditional criterion is
therefore not presented as an established solution or as an independently
verified new reduction in the literature.

% Keep the sources heading with a useful initial bibliography block.
\Needspace{8\baselineskip}
% END RESEARCH ADDITION 203
\subsection{Sources}
\begin{itemize}\raggedright
\item[\textnormal{[S1]}]\hypertarget{source:30006980:1}{} Selcuk Barlak and Xin Li, Cartan Subalgebras and the UCT Problem, arXiv1511.02697v3,21June2017,18pages.
\url{https://arxiv.org/pdf/1511.02697v3}.
{\small\textit{Catalogue formulation source.}}
% BEGIN REFERENCE ADDITION 203
\item[\textnormal{[E1]}]\hypertarget{extra:30006980:1}{}
C. Schafhauser, A. Tikuisis and S. White,
\emph{Nuclear $C^*$-algebras: 99 problems}, arXiv:2506.10902v2,
8 May 2026 (manuscript dated 11 May 2026), Section 3, especially Problem II
and the discussion of the bootstrap category.
\url{https://arxiv.org/abs/2506.10902v2}.
\item[\textnormal{[E2]}]\hypertarget{extra:30006980:2}{}
N. P. Brown, S. L. Browne, R. Willett and J. Wu,
\emph{The UCT problem for nuclear $C^*$-algebras},
arXiv:2005.03184v3 (2021), Theorems 2.1--2.2 and 4.6;
Rocky Mountain Journal of Mathematics 52 (2022), beginning at p. 817.
\url{https://arxiv.org/abs/2005.03184}.
\item[\textnormal{[E3]}]\hypertarget{extra:30006980:3}{}
R. Willett and G. Yu, \emph{Controlled $KK$-theory, and a Milnor exact
sequence}, arXiv:2011.10906, Theorems 1.1--1.2 and Section 1.2
(on Schochet's Milnor sequence and the $\lim^1$ term).
\url{https://arxiv.org/abs/2011.10906}.
\item[\textnormal{[E4]}]\hypertarget{extra:30006980:4}{}
R. Willett and G. Yu, \emph{The Universal Coefficient Theorem for
$C^*$-Algebras with Finite Complexity}, Memoirs of the European
Mathematical Society 8 (2024), 108 pp.; introduction and main UCT theorem.
DOI: 10.4171/MEMS/8.
\url{https://ems.press/books/mems/280}.
% END REFERENCE ADDITION 203
\end{itemize}

\end{document}
